\chapter{Hipotez Testlerinin Mantığı: H₀, H₁, Test İstatistiği, Anlamlılık Düzeyi ve p-Değeri}
\label{ch:hypothesis}

\begin{prismbox}[PrismLight]{Öğrenme çıktıları}
Bu bölümün sonunda:
\begin{itemize}
  \item araştırma sorusunu parametre üzerinden $H_0$ ve $H_1$ biçiminde yazabilecek,
  \item test istatistiği, anlamlılık düzeyi ve \pval-değerinin rollerini ayırabilecek,
  \item tek ve çift yönlü test arasında araştırma sorusuna göre seçim yapabilecek,
  \item bir test sonucunu güven aralığı ve bağlamla birlikte yorumlayabileceksiniz.
\end{itemize}
\end{prismbox}

\begin{prerequisitebox}
Parametre, tahmin, standart hata, güven aralığı ve örnekleme dağılımı
kavramları bilinmelidir. Standartlaştırılmış bir istatistiğin referans
dağılımla karşılaştırılması fikri hatırlanmalıdır.
\end{prerequisitebox}

Nokta tahmini ve güven aralığı için sırasıyla Bölüm~\ref{ch:point-estimation}
\ifimoders
ve Bölüm~\ref{ch:confidence-intervals} önce okunabilir.
\else
ve Bölüm~\ref{ch:confidence-intervals}; benzetim temelli test mantığı için
Bölüm~\ref{sec:randomization-test} önce okunabilir.
\fi

\section{Araştırma sorusundan hipoteze}
\label{sec:hypotheses}

Sıfır hipotezi $H_0$, veri üretim süreci hakkında sınanan temel iddiayı;
alternatif hipotez $H_1$ ise kanıt aranan karşıt durumu ifade eder.
Hipotezler veri incelenmeden önce hedef parametre üzerinden yazılmalıdır.

\section{Test istatistiği}

Test istatistiği, gözlenen tahminin $H_0$ altında beklenen değerden standart
hata birimleriyle ne kadar uzak olduğunu ölçer. Genel yapı
\[
\text{test istatistiği}=
\frac{\text{gözlenen tahmin}-\text{$H_0$ değeri}}
{\text{tahminin standart hatası}}
\]
biçimindedir.

\section{Hipotez testinin altı adımı}

\begin{enumerate}
  \item Hedef evreni, parametreyi ve araştırma sorusunu tanımlayın.
  \item $H_0$ ve $H_1$ hipotezlerini veri incelenmeden önce yazın.
  \item Anlamlılık düzeyini ve test yönünü belirleyin.
  \item Veri yapısına uygun test istatistiğini seçip varsayımları değerlendirin.
  \item Test istatistiğini ve \pval-değerini hesaplayın.
  \item Kararı etki tahmini, güven aralığı ve araştırmanın sınırlılıklarıyla raporlayın.
\end{enumerate}

Bu sıralama, test seçiminin yazılım menüsünden değil araştırma sorusu ve veri
üretim sürecinden başlamasını sağlar.

\section{Anlamlılık düzeyi ve \pval-değeri}
\label{sec:pvalue}

\begin{definitionbox}
Anlamlılık düzeyi $\alpha$, testten önce belirlenen hata eşiğidir. \pval-değeri,
$H_0$ doğru kabul edildiğinde gözlenen kadar veya daha uç bir sonuç elde etme
olasılığıdır. $p\leq\alpha$ olduğunda $H_0$ reddedilir; aksi durumda
reddetmek için yeterli kanıt bulunmadığı belirtilir.
\end{definitionbox}

\begin{mistakebox}
\pval-değeri $H_0$'ın doğru olma olasılığı, etkinin büyüklüğü veya sonucun yeniden
üretileceğinin garantisi değildir \citep{wasserstein2016}.
\end{mistakebox}

\section{Tek ve çift yönlü test}
\label{sec:test-direction}

Çift yönlü alternatif $H_1:\theta\neq\theta_0$ biçiminde iki yöndeki
sapmaları değerlendirir. Tek yönlü alternatif yalnız araştırma sorusunda
önceden gerekçelendirilmiş bir yönü sınar. Veri görüldükten sonra test yönü
değiştirilemez.

Şekil~\ref{fig:pvalue-two-tails}'de, $H_0$ altında standart normal
dağılıma sahip bir test istatistiği için $z_{\text{göz}}=2$ örneği
gösterilir. Çift yönlü testte yalnız sağ kuyruk değil, aynı uzaklıktaki
sol kuyruk da sayılır:
\[
p=\P(Z\leq-2)+\P(Z\geq2)
  =2\bigl[1-\Phi(2)\bigr]\approx0{,}0455.
\]
Bu değer, sıfır hipotezinin doğru olma olasılığı değil, $H_0$ altında
gözlenen kadar veya daha uç sonuçların toplam olasılığıdır.

\begin{figure}[H]
\centering
\begin{tikzpicture}[x=1.25cm,y=1cm,font=\footnotesize]
  \fill[PrismSubheadingBlue!35]
    (-3.8,0)--plot[domain=-3.8:-2,samples=45]
      (\x,{2.3*exp(-\x*\x/2)})--(-2,0)--cycle;
  \fill[PrismSubheadingBlue!35]
    (2,0)--plot[domain=2:3.8,samples=45]
      (\x,{2.3*exp(-\x*\x/2)})--(3.8,0)--cycle;
  \draw[PrismBlue,line width=1pt,domain=-3.8:3.8,samples=100,smooth]
    plot (\x,{2.3*exp(-\x*\x/2)});
  \draw[->,PrismGray] (-3.9,0)--(3.95,0) node[right] {$z$};
  \foreach \x in {-2,2}
    \draw[PrismBlue,dashed] (\x,0)--(\x,0.8);
  \node[below] at (0,0) {$0$};
  \node[below] at (-2,0) {$-|z_{\text{göz}}|=-2$};
  \node[below] at (2,0) {$z_{\text{göz}}=2$};
  \node[align=center] at (0,2.7) {$H_0$ altındaki referans dağılım};
  \node at (-2.85,1.15) {Sol kuyruk: $p/2$};
  \node at (2.85,1.15) {Sağ kuyruk: $p/2$};
  \draw[->,PrismBlue] (-2.85,0.9)--(-2.5,0.07);
  \draw[->,PrismBlue] (2.85,0.9)--(2.5,0.07);
\end{tikzpicture}
\caption{Standart normal referans altında çift yönlü \pval-değeri.
İki gölgeli kuyruğun temsil ettiği toplam olasılık, $z_{\text{göz}}=2$ için yaklaşık
$0{,}0455$'tir. Dikey ölçek görsel açıklık için büyütülmüştür.}
\label{fig:pvalue-two-tails}
\end{figure}

\section{Kritik bölge ve \pval-değeri yaklaşımı}

Hipotez testi aynı kararı iki eşdeğer yolla verebilir. Standart normal
referans altında çift yönlü bir testte
\[
|Z|\geq z_{1-\alpha/2}
\]
olduğunda gözlenen değer kritik bölgededir. Üst yönlü bir testte ise
$Z\geq z_{1-\alpha}$ koşulu kullanılır. Alternatif olarak \pval-değeri doğrudan
$\alpha$ ile karşılaştırılır. Aynı model, yön ve anlamlılık düzeyi
kullanıldığında iki yaklaşımın kararı aynı olmalıdır.

Kritik değer yaklaşımı karar sınırını, \pval-değeri yaklaşımı ise gözlenen
sonucun sınırın ne kadar ötesinde veya gerisinde olduğunu gösterir. Her iki
yaklaşımda da etki tahmini ve güven aralığı ayrıca raporlanmalıdır.

\section{Sonucun yorumlanması}
\label{sec:test-reporting}

Bir test sonucu; test istatistiği, serbestlik derecesi, \pval-değeri, tahmin,
güven aralığı ve araştırma bağlamıyla birlikte raporlanmalıdır. ``Anlamlı''
ifadesi tek başına uygulamalı önemi göstermez.

\section{\pval-değeri bir kanıt ölçüsüdür, kararın tamamı değildir}

\pval-değeri küçüldükçe gözlenen sonucun $H_0$ altında açıklanması güçleşir;
ancak 0,049 ile 0,051 arasında bilimsel bakımdan keskin bir sınır yoktur.
Örneklem hacmi çok büyük olduğunda önemsiz bir fark küçük \pval-değeri üretebilir;
örneklem küçük olduğunda ise uygulamada önemli bir fark belirsizlik nedeniyle
istatistiksel eşiği geçemeyebilir. Eşik merkezli mekanik yorumların yerine
tahmin, belirsizlik, bağlam ve çalışma tasarımının birlikte değerlendirilmesi
önerilmektedir \citep{wasserstein2019}.

Çift yönlü bir testte yüzde 95 güven aralığı $H_0$ değerini dışlıyorsa aynı
veriye dayanan test genellikle $\alpha=0{,}05$ düzeyinde anlamlıdır. Güven
aralığı ayrıca etkinin yönünü, olası büyüklüklerini ve tahmin hassasiyetini
gösterdiği için \pval-değerinden daha fazla bilgi taşır.

\section{Çözümlü düşünme örneği}

Bir öğretim programının ortalama başarı puanını 50'den değiştirdiği iddia
edilsin. Hipotezler $H_0:\mu=50$ ve $H_1:\mu\neq50$ olarak kurulur. Analiz
sonucunda $t=2{,}50$ ve $p=0{,}020$ elde edilirse, $\alpha=0{,}05$ düzeyinde
$H_0$ reddedilir. Bu karar ``$H_0$ kesinlikle yanlıştır'' anlamına gelmez;
verinin $H_0$ altında beklenenden yeterince sıra dışı olduğunu belirtir.

Aynı sonuç için $\alpha=0{,}01$ önceden seçilmiş olsaydı $H_0$ reddedilemezdi.
Veri değişmemiş, yalnız karar eşiği değişmiştir. Bu nedenle anlamlılık düzeyi
sonuca bakılmadan önce belirlenmelidir.

\section{Çözümlü uygulama: akademik destek programı}

Bir üniversitede önceki yılların ortalama akademik uyum puanı 70'tir. Yeni bir
destek programına katılan $n=50$ öğrencinin ortalaması 72,1 ve standart hatası
1 puan bulunmuştur. Programın puanı herhangi bir yönde değiştirip
değiştirmediği araştırılsın.

\begin{enumerate}
  \item \textbf{Parametre:} Programı alan hedef evrendeki ortalama uyum puanı $\mu$.
  \item \textbf{Hipotezler:} $H_0:\mu=70$ ve $H_1:\mu\neq70$.
  \item \textbf{Test istatistiği:}
  \[
  t=\frac{72{,}1-70}{1}=2{,}10,\qquad sd=49.
  \]
  \item \textbf{Kanıt:} Çift yönlü \pval-değeri yaklaşık $0{,}041$'dir.
  \item \textbf{Belirsizlik:} Ortalama farkının yüzde 95 güven aralığı
  $[0{,}09;4{,}11]$ puandır.
  \item \textbf{Karar ve yorum:} $\alpha=0{,}05$ düzeyinde $H_0$ reddedilir.
  Program grubunun ortalaması referans değerden yüksek görünmektedir; ancak
  farkın küçük olabileceği ve karşılaştırma grubu yoksa sonucun doğrudan
  programa bağlanamayacağı belirtilmelidir.
\end{enumerate}

\begin{researchbox}
Referans değer başka yıllardan geliyorsa öğrenci profili, ölçme aracı veya
uygulama koşulları değişmiş olabilir. İstatistiksel test bu tasarım farklarını
kendiliğinden denetlemez.
\end{researchbox}

\begin{outputguidebox}
Yazılım çıktısında önce sınanan hipotezin ve test yönünün doğru tanımlandığını
kontrol edin. Ardından test istatistiği, serbestlik derecesi, \pval-değeri, etki
tahmini ve güven aralığını birlikte okuyun. Yazılımın verdiği \pval-değeri
araştırma tasarımındaki yanlılıkları düzeltemez.
\end{outputguidebox}

\section{Yazılım uygulaması: test sonucu ve güven aralığı}

\begin{softwarebox}
\begin{lstlisting}
from scipy import stats

t_value, df = 2.10, 49
p_value = 2 * stats.t.sf(abs(t_value), df=df)
critical = stats.t.ppf(.975, df=df)
diff_ci = (2.1 - critical, 2.1 + critical)

print(p_value)
print(diff_ci)
\end{lstlisting}
Sonuç yaklaşık $p=0{,}041$ ve fark için $[0{,}09;4{,}11]$ aralığını verir.
Test ile güven aralığının aynı veri, model ve yön tanımını kullandığı
doğrulanmalıdır.
\end{softwarebox}

\begin{assumptionbox}
\pval-değerinin geçerliliği yalnız hesap formülüne değil, örnekleme tasarımına,
bağımsızlığa, test yönünün önceden belirlenmesine ve kullanılan modelin
uygunluğuna bağlıdır.
\end{assumptionbox}

\section{Çözümlü problemler}

\subsection*{Problem 1: Araştırma sorusundan hipotezlere}

Bir destek programının öğrencilerin ortalama öz yeterlik puanını bilinen 50
puanlık düzeyin üzerine çıkardığı iddia edilmektedir. Hedef parametreyi,
hipotezleri ve uygun test yönünü belirleyin.

\textbf{Çözüm.} Hedef parametre, programın uygulanacağı evrendeki ortalama öz
yeterlik puanı $\mu$'dür. Araştırma iddiası yalnız artış yönünde ve veri
görülmeden önce kurulmuşsa
\[
H_0:\mu\leq50,
\qquad
H_1:\mu>50
\]
biçiminde üst yönlü test uygundur. Test istatistiğinin sıfır hipotezi altındaki
dağılımı sınır değer olan $\mu=50$ üzerinden hesaplanır.

Programın puanı herhangi bir yönde değiştirip değiştirmediği sorulsaydı
$H_1:\mu\neq50$ biçiminde çift yönlü test gerekirdi. Yalnız gözlenen ortalama
50'nin üzerinde çıktığı için sonradan tek yönlü teste geçilemez.

\textbf{Raporlama örneği.} ``Programın ortalama öz yeterlik puanını 50'nin
üzerine çıkaracağı yönündeki önceden belirlenmiş hipotez üst yönlü testle
incelendi.''

\subsection*{Problem 2: Aynı test istatistiğinde tek ve çift yönlü \pval-değeri}

Bir araştırmada gözlenen ortalama 74, sıfır hipotezindeki değer 70 ve
tahminin standart hatası 2'dir. Büyük örneklem normal yaklaşımı altında test
istatistiğini, üst yönlü ve çift yönlü \pval-değerlerini bulun.

\textbf{Çözüm.} Test istatistiği
\[
Z=\frac{74-70}{2}=2
\]
olur. Standart normal tabloda $\Phi(2)=0{,}9772$ olduğundan üst yönlü
\pval-değeri
\[
p_{\text{üst}}=1-0{,}9772=0{,}0228
\]
ve çift yönlü \pval-değeri
\[
p_{\text{çift}}=2(0{,}0228)=0{,}0456
\]
olarak bulunur. Her ikisi de 0,05'ten küçüktür; ancak tek yönlü \pval-değeri daha
küçüktür çünkü karşı kuyruğu değerlendirmez.

Bu fark, daha küçük \pval-değeri elde etmek için tek yönlü test seçmeyi haklı
çıkarmaz. Test yönü veri görülmeden önce araştırma sorusuna göre belirlenir.

\subsection*{Problem 3: Kritik değerle karar verme}

Çift yönlü standart normal testte gözlenen değer $Z=1{,}85$ olsun. Sonucu
$\alpha=0{,}05$ ve $\alpha=0{,}10$ düzeylerinde kritik değer yaklaşımıyla
değerlendirin. Çift yönlü \pval-değeri yaklaşık 0,064'tür.

\textbf{Çözüm.} Yüzde 5 düzeyinde kritik değerler yaklaşık $\pm1{,}96$'dır.
\[
|1{,}85|<1{,}96
\]
olduğu için $H_0$ reddedilemez. Aynı karar $p=0{,}064>0{,}05$
karşılaştırmasından da elde edilir.

Yüzde 10 düzeyinde kritik değerler yaklaşık $\pm1{,}645$'tir.
\[
|1{,}85|>1{,}645
\]
olduğundan $H_0$ reddedilir; aynı biçimde $0{,}064<0{,}10$'dur. Veri ve etki
tahmini değişmemiş, yalnız önceden kabul edilen Tip I hata eşiği değişmiştir.
Bu nedenle sonuç ``yüzde 10'da anlamlı, yüzde 5'te anlamlı değil'' biçiminde
verilebilir; kanıtın niteliği yapay olarak iki bütünüyle farklı sınıfa
ayrılmamalıdır.

\subsection*{Problem 4: Örneklem hacmi ve uygulamalı önem}

İki çalışmada aynı 2 puanlık ortalama fark gözlenmiştir. İlk çalışmada
standart hata 1, ikinci çalışmada 0,25'tir. Sıfır fark hipotezi için test
istatistiklerini hesaplayın. Uygulamada en az 5 puanlık fark önemliyse
sonuçları yorumlayın.

\textbf{Çözüm.} İlk çalışmada
\[
Z_1=\frac{2}{1}=2
\]
ve çift yönlü \pval-değeri yaklaşık 0,046'dır. İkinci çalışmada
\[
Z_2=\frac{2}{0{,}25}=8
\]
olur ve \pval-değeri son derece küçüktür. Daha düşük standart hata, aynı etki
tahminini sıfırdan çok daha kesin biçimde ayırmıştır.

Bununla birlikte iki çalışmada da gözlenen fark 2 puandır ve önceden
belirlenen 5 puanlık uygulamalı önem eşiğinin altındadır. İkinci sonuç çok
güçlü istatistiksel kanıt sunabilir fakat etkinin uygulamadaki büyüklüğü yine
küçük olabilir. \pval-değeri etki büyüklüğünün ölçüsü değildir.

\subsection*{Problem 5: Test ile güven aralığının birlikte okunması}

Bir öğretim uygulamasının ortalama etkisi 3 puan ve yüzde 95 güven aralığı
$[0{,}5;5{,}5]$ bulunmuştur. Sıfır fark hipotezini ve en az 4 puanlık
uygulamalı önem eşiğini dikkate alarak sonucu yorumlayın.

\textbf{Çözüm.} Aralık sıfırı dışladığı için aynı çift yönlü test
$\alpha=0{,}05$ düzeyinde $H_0$'ı reddeder. Etkinin pozitif olduğuna ilişkin
istatistiksel kanıt vardır. Ancak güven aralığı 0,5 ile 5,5 arasında geniş bir
etki kümesini kapsar.

Aralığın bir bölümü 4 puanlık uygulamalı önem eşiğinin altında, bir bölümü
üzerindedir. Dolayısıyla etkinin yönü konusunda kanıt bulunmasına rağmen
uygulamada önemli büyüklüğe ulaşıp ulaşmadığı kesin değildir. Yalnız
``istatistiksel olarak anlamlıdır'' demek belirsizliğin bu yönünü gizler.

\subsection*{Problem 6: Oran için hipotez testi}

Bir dijital platformda öğrencilerin yarısının bir etkinliği tamamladığı
$H_0:p=0{,}50$ hipotezi sınanacaktır. Rastgele seçilen 200 öğrencinin yüzde
57'si etkinliği tamamlamıştır. Büyük örneklem yaklaşımıyla test istatistiğini
ve çift yönlü \pval-değerini yaklaşık olarak bulun.

\textbf{Çözüm.} Sıfır hipotezi altında oran standart hatası
\[
SE_0=\sqrt{\frac{0{,}50(1-0{,}50)}{200}}
=\sqrt{0{,}00125}\approx0{,}03536
\]
olur. Test istatistiği
\[
Z=\frac{0{,}57-0{,}50}{0{,}03536}\approx1{,}98
\]
ve çift yönlü \pval-değeri yaklaşık 0,048'dir. Yüzde 5 düzeyinde $H_0$ sınırda
reddedilir; veri, tamamlanma oranının yüzde 50'den farklı olduğuna ilişkin
ılımlı kanıt sunar.

Test hesabında standart hata $H_0$ değeri olan 0,50 ile kurulmuştur. Güven
aralığında ise standart hata genellikle gözlenen $\hat p=0{,}57$ üzerinden
tahmin edilir. İki hesap yakın olsa da amaçları aynı değildir.

\begin{mistakebox}
$p=0{,}048$ sonucu, sıfır hipotezinin yanlış olma olasılığının yüzde 95,2
olduğu anlamına gelmez. Bu değer, $H_0$ altında gözlenen kadar veya daha uç
bir oran farkının olasılığını özetler.
\end{mistakebox}

\begin{outputguidebox}
Bir test çıktısında araştırma hipotezinin yönünü, sıfır hipotezindeki değeri,
test istatistiğini, serbestlik derecesini ve tek/çift yönlü \pval-değerini kontrol
edin. Ardından etki tahmini ile güven aralığını okuyun. ``Sig.'' sütununu tek
başına raporlamak araştırma sonucunu tamamlamaz.
\end{outputguidebox}

\section{Literatür verisiyle yazılım uygulamaları}
\label{sec:spss-b09}

\textbf{Amaç ve veri.} \texttt{b09.csv} ile öğretim amaçlı
$H_0:\mu=10$ ve $H_1:\mu\ne10$ hipotezlerini sınayın
\citep{cortez2008data}. 10 puan, burada seçilen referans değeridir.
İki yönlü tek örneklem $t$ testinde anlamlılık düzeyi $\alpha=0{,}05$,
güven düzeyi yüzde 95 olarak belirlenmiştir.


\textbf{Ortak veri.} Doğrulanan B09 uygulamasında üç yazılım da
\nolinkurl{companion/spss/csv/b09.csv} dosyasını kullanmıştır. Dosya
395 satır ve 10 sütun içerir; kayıtların tamamı analiz edilir.
B06'nın seçim göstergeleri ve ağırlığı burada kullanılmaz.
Bu karşılaştırma Excel içe aktarmanın doğrulandığı anlamına gelmez.
Kodları \texttt{companion} klasörünü içeren paket kökünden çalıştırın;
veri sözlüğü ve hazırlık için bkz. sayfa~\pageref{sec:spss-guide}.

\subsection{IBM SPSS uygulaması}
\textbf{Başlamadan önce.} Aşağıdaki komut dosyası B09 CSV verisini
açar ve analiz eder. Menü yolunu kullanacaksanız önce veri açma ve
tanımlama komutlarını \texttt{EXECUTE} satırı dahil çalıştırın.
B08'den kalan tek satırlık özet yerine ham veriyi yeniden açın;
filtre, ağırlık ve dosya bölme kapalı olmalıdır. \texttt{GET DATA}
dosya açma hatası verirse önce \texttt{/FILE} yolunu düzeltin ve
kodu baştan çalıştırın. Veri açılmadan sonraki komutlar çalışmaz.

\begin{spssadimlar}
\spssadim{Önce \emph{Analyze $\to$ Descriptive Statistics $\to$ Explore} ile \texttt{g3} histogramını inceleyin. Referans değerini sonuçlara bakarak değiştirmeyin.}
\spssadim{\emph{Analyze $\to$ Compare Means $\to$ One-Sample T Test} yolunu açın. Yeni sürümlerde orta menü \emph{Compare Means and Proportions} olarak görünebilir.}
\spssadim{\texttt{g3} değişkenini \emph{Test Variable(s)} alanına taşıyın; \emph{Test Value} alanına 10 yazın.}
\spssadim{\emph{Options} altında güven düzeyini yüzde 95 bırakın; \emph{Continue $\to$ OK} seçin. Bu uygulamanın alternatif hipotezi iki yönlüdür.}
\spssadim{\emph{One-Sample Test} tablosunda \emph{t}, \emph{df} ve iki yönlü anlamlılık sütununu okuyun. Tek yönlü anlamlılık sütunu da varsa bu uygulamada onu kullanmayın.}
\end{spssadimlar}

{\raggedright
\noindent\textbf{Komutların görevi.} \texttt{TESTVAL=10} sıfır hipotezindeki ortalamayı, \texttt{VARIABLES=g3} test değişkenini belirtir. \texttt{CI(.95)} fark için yüzde 95 güven aralığı ister.\par}

\textbf{Uygulama kodu:} \texttt{b09}. Doğrulanan veri dosyası
\texttt{companion/spss} altındaki \texttt{csv/b09.csv} dosyasıdır.

\begin{lstlisting}[style=prismSPSS]
SET DECIMAL=DOT.
GET DATA /TYPE=TXT
 /FILE='companion/spss/csv/b09.csv'
 /ENCODING='UTF8'
 /ARRANGEMENT=DELIMITED /DELCASE=LINE
 /DELIMITERS="," /QUALIFIER='"'
 /FIRSTCASE=2
 /VARIABLES=id F8.0 school F8.0 sex F8.0 age F8.0
 studytime F8.0 absences F8.0 g1 F8.0 g2 F8.0
 g3 F8.0 pass10 F8.0.
FILTER OFF.
WEIGHT OFF.
SPLIT FILE OFF.
VARIABLE LABELS
 g3 'Yil sonu matematik notu'
 pass10 'Yil sonu notu 10 veya uzerinde'.
VALUE LABELS school 1 'GP' 2 'MS'.
VALUE LABELS pass10
 0 '10 puanin altinda' 1 '10 puan ve uzerinde'.
VARIABLE LEVEL
 id school sex pass10 (NOMINAL)
 studytime (ORDINAL)
 age absences g1 g2 g3 (SCALE).
FORMATS g3 (F14.9).
EXECUTE.
EXAMINE VARIABLES=g3
 /PLOT=HISTOGRAM /STATISTICS=DESCRIPTIVES
 /CINTERVAL=95.
T-TEST /TESTVAL=10 /VARIABLES=g3
 /CRITERIA=CI(.95).
\end{lstlisting}

\begin{outputguidebox}
\textbf{Paylaşılan SPSS çıktısıyla doğrulanan değerler.}
\emph{One-Sample Statistics} tablosunda $n=395$, ortalama \SPMatMean\ ve
standart hata \SPMatSE\ bulunmuştur. \emph{One-Sample Test} tablosunda
$t(394)=1{,}801$, çift yönlü $p=0{,}072$ gösterilmiştir. Yüzde 5 düzeyinde
$H_0$ reddedilmez; bu, ortalamanın kesinlikle 10 olduğunun veya farkın
önemsiz olduğunun kanıtı değildir. Test tablosundaki güven aralığı
\emph{ortalama farkı} içindir: Explore çıktısındaki ortalama aralığının
iki sınırından da 10 çıkarılır. Histogram, sıfır yığılmasını değerlendirmek
içindir; büyük örneklem bağımsızlık sorununu çözmez.
\end{outputguidebox}


\par\noindent\begin{minipage}{\linewidth}
\subsection{R uygulaması}
\label{sec:r-gercek-b09}

\texttt{mu=10} sıfır hipotezini belirler. \texttt{conf.int} ortalama için aralıktır; 10 çıkarılınca fark aralığı elde edilir.

\begin{lstlisting}[style=prismR]
veri <- read.csv("companion/spss/csv/b09.csv")
stopifnot(nrow(veri) == 395, ncol(veri) == 10,
          !anyNA(veri),
          all(veri$pass10 == as.integer(veri$g3 >= 10)))
sonuc <- t.test(veri$g3, mu=10,
               alternative="two.sided", conf.level=.95)
ozet <- c(
  n=nrow(veri), ort=mean(veri$g3), ss=sd(veri$g3),
  sh_ort=sd(veri$g3)/sqrt(nrow(veri)),
  t=unname(sonuc$statistic), df=unname(sonuc$parameter),
  p_iki_yonlu=sonuc$p.value, fark=mean(veri$g3)-10,
  ort_alt=sonuc$conf.int[1], ort_ust=sonuc$conf.int[2],
  fark_alt=sonuc$conf.int[1]-10,
  fark_ust=sonuc$conf.int[2]-10)
print(sonuc)
print(ozet, digits=15)
print(sum(veri$g3 == 0))
\end{lstlisting}

\textbf{Çıktıyı okuma.} Paylaşılan R çıktısı 395 satır, 10 sütun ve
her sütunda sıfır eksik değer gösterir. Test çıktısındaki
\texttt{p-value}, burada iki yönlü olasılıktır. Ayrıntılı özette
\texttt{ort\_alt} ve \texttt{ort\_ust} ortalama için;
\texttt{fark\_alt} ve \texttt{fark\_ust} ise ortalama eksi 10 için
sınırlardır. Sıfır notlu 38 kayıt analizde tutulmuştur.
\end{minipage}\par

\par\noindent\begin{minipage}{\linewidth}
\subsection{Python uygulaması}
\label{sec:python-b09}

\texttt{sonuc} iki yönlü testin istatistiğini ve $p$ değerini tutar.
\texttt{ddof=1} örneklem standart sapmasını seçer; ortalama aralığından
10 çıkarılarak SPSS fark aralığıyla karşılaştırılır.

\begin{lstlisting}[style=prismPythonApp]
import pandas as pd
import numpy as np
from scipy import stats

veri = pd.read_csv("companion/spss/csv/b09.csv")
assert veri.shape == (395, 10)
assert not veri.isna().any().any()
assert veri.pass10.eq((veri.g3 >= 10).astype(int)).all()
gozlem_sayisi = len(veri)
ortalama = veri.g3.mean()
standart_sapma = veri.g3.std(ddof=1)
standart_hata = standart_sapma / np.sqrt(gozlem_sayisi)
sonuc = stats.ttest_1samp(
    veri.g3, popmean=10, alternative="two-sided")
ortalama_alt, ortalama_ust = stats.t.interval(
    .95, df=gozlem_sayisi-1,
    loc=ortalama, scale=standart_hata)
ozet = pd.Series({
    "n": gozlem_sayisi, "ort": ortalama, "ss": standart_sapma,
    "sh_ort": standart_hata, "t": sonuc.statistic,
    "df": gozlem_sayisi-1, "p_iki_yonlu": sonuc.pvalue,
    "fark": ortalama-10,
    "ort_alt": ortalama_alt, "ort_ust": ortalama_ust,
    "fark_alt": ortalama_alt-10, "fark_ust": ortalama_ust-10
})
print(ozet.to_string(float_format=lambda deger: f"{deger:.13f}"))
print(int(veri.g3.eq(0).sum()))
\end{lstlisting}

\textbf{Çıktıyı okuma.} Paylaşılan Python çıktısında da 395 satır,
10 sütun, sıfır eksik değer ve 38 sıfır notlu kayıt vardır.
On iki özet değer, R sonuçları 13 ondalık basamağa yuvarlandığında
eşleşir. Her iki yazılımda ana veri korunur. Buradaki karşılaştırma
sayısal çıktılara aittir; R ve Python histogramları doğrulanmamıştır.
\end{minipage}\par

\subsection{Sonuçların karşılaştırılması ve ortak rapor}
12 Eylül 2026'da paylaşılan SPSS tabloları ve histogramı ile R ve
Python metin çıktıları incelenmiştir. Gözlem sayısı, ortalama,
standart sapma, standart hata, $t$, serbestlik derecesi, iki yönlü
$p$, ortalama farkı ve iki aralığın dört sınırından oluşan on iki
özet değer R ve Python'da 13 ondalık basamağa yuvarlandığında
eşleşmiştir. SPSS'in ilgili değerleri gösterilen yuvarlama
hassasiyetlerinde uyumludur. Proje CSV'sinden yapılan hesaplar
sonuçları destekler. Aşağıdaki tablolar çıktıların kitap düzenine
uyarlanmış özetidir; ekran görüntüsü değildir.

\begin{table}[H]
\centering
\caption{B09'da yıl sonu notunun doğrulanmış betimsel özeti}
\label{tab:b09-dogrulanmis-betimsel}
\begin{tabular}{@{}lr@{}}
\toprule
Ölçü & Değer \\
\midrule
Gözlem sayısı & 395 \\
Ortalama & 10,415190 \\
Standart sapma & 4,581443 \\
Ortalamanın standart hatası & 0,230517 \\
\bottomrule
\end{tabular}
\par\smallskip
{\footnotesize Ondalıklı değerler altı basamağa yuvarlanmıştır.
Standart sapma 394 böleniyle, standart hata $s/\sqrt{395}$ ile
hesaplanmıştır. \texttt{g3} için eksik değer yoktur; sıfır notlu
38 kayıt analizde korunmuştur.\par}
\end{table}

\begin{table}[H]
\centering
\caption{B09'da 10 puan referansına karşı iki yönlü tek örneklem testi}
\label{tab:b09-dogrulanmis-test}
\begin{tabular}{@{}lr@{}}
\toprule
Ölçü & Değer \\
\midrule
$t$ istatistiği & 1,801122 \\
Serbestlik derecesi & 394 \\
İki yönlü $p$ değeri & 0,072448 \\
Ortalama farkı ($\bar x-10$) & 0,415190 \\
Ortalamanın \%95 güven aralığı: alt & 9,961992 \\
Ortalamanın \%95 güven aralığı: üst & 10,868388 \\
Farkın \%95 güven aralığı: alt & $-0{,}038008$ \\
Farkın \%95 güven aralığı: üst & 0,868388 \\
\bottomrule
\end{tabular}
\par\smallskip
{\footnotesize Ondalıklı değerler altı basamağa yuvarlanmıştır;
ara hesaplarda yuvarlanmamış değerler kullanılmıştır. SPSS test
tablosu farkın aralığını, Explore tablosu ortalamanın aralığını
verir. Fark aralığının iki sınırına da 10 eklenince ortalama
aralığı elde edilir.\par}
\end{table}

\Needspace{15\baselineskip}
\begin{outputguidebox}
Tablo~\ref{tab:b09-dogrulanmis-test} için $p=0{,}072448>0{,}05$
olduğundan $H_0:\mu=10$ reddedilmez. Bu karar ortalamanın kesinlikle
10 olduğunu, iki değerin eşdeğerliğini veya farkın uygulamada
önemsizliğini kanıtlamaz. Ortalama aralığı 10'u, fark aralığı 0'ı
içerir; aynı iki yönlü test ve güven düzeyi için bu sonuçlar uyumludur.

Paylaşılan SPSS tablosunda \emph{One-Sided p} sütunu 0,036,
\emph{Two-Sided p} sütunu 0,072 gösterir. B09'un alternatif hipotezi
iki yönlü olduğundan ikinci sütun kullanılır. Yardım açıklamasındaki
tek yönlü anlamlılık vurgusu, bu uygulamanın kararını değiştirmez.
Sonuçlara bakarak testin yönünü değiştirmeyin; renk yerine hipotezi
ve ilgili sütunun başlığını esas alın.

SPSS histogramında sıfırda ayrı bir yığılma vardır; bu 38 kayıt
gerekçesiz biçimde silinmemiştir. Histogram bağımsızlığı veya temsil
gücünü doğrulamaz. Hesaplar öğretim amacıyla kullanılan bağımsız
gözlemler modeline dayanır; iki okuldan gelen verinin daha geniş bir
öğrenci evrenine genellenebilirliği, yazılımlar aynı sonucu verdiği
için kanıtlanmış olmaz.
\end{outputguidebox}

\Needspace{10\baselineskip}
\textbf{Raporlama örneği (doğrulanmış çıktılarla).}
``395 kayıtta yıl sonu notunun ortalaması 10,415190, standart sapması
4,581443 bulunmuştur. Öğretim amaçlı 10 puan referansına karşı iki
yönlü tek örneklem $t$ testinde yüzde 5 düzeyinde sıfır hipotezi
reddedilmemiştir ($t(394)=1{,}801122$, $p=0{,}072448$).
Ortalama farkı 0,415190 puan, farkın yüzde 95 güven aralığı
$[-0{,}038008;0{,}868388]$ bulunmuştur. Ortalamanın yüzde 95 güven
aralığı $[9{,}961992;10{,}868388]$ olarak hesaplanmıştır. Sıfır notlu
38 kayıt korunmuştur. Bağımsız gözlemler modeline dayanan bu sonuçlar,
iki okulla sınırlı veri kaynağının daha geniş bir öğrenci evrenine
genellenebilirliğini kanıtlamamaktadır.''


\textbf{Görev.} Farkın güven aralığına 10 ekleyerek ortalamanın
aralığını doğrulayın. Neden 0,036 yerine 0,072 sütununun kullanıldığını
hipotezlerle açıklayın. $H_0$'ın reddedilmemesi ile eşdeğerlik
kanıtı arasındaki ayrımı ve bağımsızlık varsayımının histogramdan
doğrulanamayacağını belirtin.

\section{R ile hipotez testi}
\label{sec:r-b09}

\textbf{Soru.} Bölümdeki çift yönlü testte gözlenen ortalama farkı 2,1,
standart hata 1 ve serbestlik derecesi 49 olsun. $H_0:\delta=0$ için
\pval-değerini, yüzde 95 güven aralığını ve $\alpha=0{,}05$ kararını bulun.

\begin{lstlisting}[style=prismR]
fark <- 2.1
standart_hata <- 1
serbestlik <- 49
alfa <- 0.05
t_degeri <- fark / standart_hata
p_degeri <- 2 * pt(abs(t_degeri), df = serbestlik,
                   lower.tail = FALSE)
kritik <- qt(1 - alfa / 2, df = serbestlik)
aralik <- fark + c(-1, 1) * kritik * standart_hata
c(t = t_degeri, p = p_degeri, kritik = kritik)
aralik
if (p_degeri < alfa) {
  print("H0 reddedilir")
} else {
  print("H0 reddedilemez")
}
pt(t_degeri, df = serbestlik, lower.tail = FALSE)
\end{lstlisting}

\begin{outputguidebox}
\textbf{Çözüm ve kontrol değerleri.} $t=2{,}10$, çift yönlü
$p=0{,}04090$ ve kritik değer 2,0096'dır. Güven aralığı
$[0{,}0904;4{,}1096]$ sıfırı içermez; yüzde 5 düzeyinde $H_0$ reddedilir.
Son satır, önceden belirlenmiş $H_1:\delta>0$ için tek yönlü
$p=0{,}02045$ verir.
\end{outputguidebox}

\textbf{Yorum.} Çift yönlü değer, mutlak test istatistiğinin üst kuyruk
olasılığının iki katıdır. Sonuca bakıp yön seçmek uygun değildir. Küçük
\pval-değeri $H_0$'ın doğru olma olasılığı veya etkinin önemli olduğu anlamına gelmez.

\textbf{Pekiştirme ve yanıt.} Önceden $\alpha=0{,}01$ seçilseydi aynı
$p=0{,}04090$ ile $H_0$ reddedilemezdi. Veri değişmez; karar eşiği değişir.

\section{Bölüm özeti}

\begin{itemize}
  \item Hipotezler veri görülmeden önce hedef parametre üzerinden kurulur.
  \item \pval-değeri $H_0$ altında verinin ne kadar sıra dışı olduğunu ölçer; $H_0$'ın olasılığı değildir.
  \item Reddedememek eşitliği veya etkisizliği kanıtlamak anlamına gelmez.
  \item Karar; tahmin, güven aralığı, etki büyüklüğü ve tasarım sınırlılıklarıyla raporlanır.
  \item Kritik bölge ve \pval-değeri yaklaşımları aynı model ve anlamlılık düzeyi
  altında aynı test kararını verir.
  \item Büyük örneklem küçük etkileri istatistiksel olarak ayırt edebilir;
  uygulamalı önem ayrıca değerlendirilmelidir.
\end{itemize}

\section*{Alıştırmalar}

\begin{enumerate}
  \item $H_0$ ve $H_1$ için bir eğitim araştırması örneği yazın.
  \item \pval-değerinin ne olduğunu ve ne olmadığını açıklayın.
  \item $H_0$ reddedilemediğinde hangi ifade kullanılmalıdır?
  \item Aynı \pval-değeri için $\alpha=0{,}05$ ve $\alpha=0{,}01$ kararlarını karşılaştırın.
  \item İstatistiksel anlamlılık ile uygulamalı önemin neden ayrı değerlendirilmesi gerektiğini açıklayın.
  \item Akademik destek örneğinde güven aralığının alt sınırını uygulamalı önem bakımından yorumlayın.
  \item $p=0{,}049$ ile $p=0{,}051$ sonuçlarının neden bütünüyle farklı kanıtlar olmadığını açıklayın.
  \item Araştırma sorusuna uygun bir tek yönlü hipotez yazın ve yön seçimini gerekçelendirin.
  \item Bir öğretim yönteminin ortalamayı azalttığı iddiası için hedef
  parametreyi ve tek yönlü hipotezleri yazın.
  \item Tahmin 62, sıfır hipotezi değeri 58 ve standart hata 2 ise test
  istatistiğini, üst yönlü ve çift yönlü yaklaşık \pval-değerlerini bulun.
  \item $Z=-2{,}20$ sonucunu yüzde 5 çift yönlü testte kritik değer
  yaklaşımıyla değerlendirin.
  \item $p=0{,}032$ sonucunda $\alpha=0{,}05$ ve $\alpha=0{,}01$ için
  kararları ve uygun sözel ifadeleri yazın.
  \item Etki tahmini 1,5 puan ve yüzde 95 güven aralığı $[0{,}2;2{,}8]$ iken
  uygulamalı önem sınırı 4 puansa sonucu yorumlayın.
  \item Fark tahmini 6 ve güven aralığı $[-1;13]$ olan bir araştırmada
  $H_0$'ın reddedilememesinin neden ``etki yoktur'' demek olmadığını açıklayın.
  \item $H_0:p=0{,}40$, $n=400$ ve $\hat p=0{,}45$ için sıfır hipotezi
  altındaki standart hatayı ve test istatistiğini hesaplayın.
  \item Aynı etki tahmininin farklı örneklem hacimlerinde farklı \pval-değerleri
  üretmesini standart hata üzerinden açıklayın.
\end{enumerate}